% Round 18. New-sectors branch. Uses the frozen GP17 arithmetic pencil.
\section{Full-height fourfold families with a deep degeneration}
\label{r18gf:section}

The local hyperplane theorem can be inserted into the arithmetic
pencils already proved in GP17.  A congruence modulo $q^2$ fixes the
new local model while keeping that frozen pencil theorem applicable
at all other places.  The extension statement below is deliberately
restricted to fields unramified at $q$; new cycles over those fields
are treated by applying the local theorems anew.

\begin{theorem}[An infinite full-height family with $N^4\ne0$]
\label{r18gf:global}
For every integer $d\ge6$ there exist a rational prime $q$, an
integral degree-$d$ form $F_1$, and an infinite arithmetic progression
of integers $a$, with finitely many excluded values, such that
\[
 X_a=\{dF_0+aF_1=0\}\subset\mathbf P^5_{\mathbf Q},
 \qquad F_0=X_0^d+\cdots+X_5^d,
 \tag{GF18.1}\label{r18gf:family}
\]
is smooth and geometrically integral and has the following properties.
\begin{enumerate}
\item At every finite place of $\mathbf Q$, every specified
geometrically homologically trivial rational cycle has ordinary
actual strong G, with its exact support and all original actual
compact tests.
\item At $q$ the variety has a strict semistable model of the type
in Theorem~\ref{r18hp:model}, and $N^4\ne0$ on $H^4$.
\item The full local statement holds for every number field $F$
which is unramified at $q$, at every finite place of $F$, for all
cycles over the local completion, including those not descending
to $\mathbf Q_q$ or to $\mathbf Q$.
\item For every such $F$ and every $1\le r\le4$, the frozen PH16
bridge supplies the real Beilinson--Bloch height on the full groups
\[
 \operatorname{CH}^r((X_a)_F)_{\mathbf Q,\mathrm{geom\,hom}}
 \ \times\
 \operatorname{CH}^{5-r}((X_a)_F)_{\mathbf Q,\mathrm{geom\,hom}}
                         \longrightarrow\mathbf R.
 \tag{GF18.2}\label{r18gf:height}
\]
\end{enumerate}
One may require $q$ to avoid any prescribed finite set of rational
primes.  In particular, for any fixed number field $F$, one may
construct a family satisfying these conclusions over $F$ by also
requiring $q\nmid\operatorname{disc}(F)$.
The theorem asserts infinitely many parameter members, not pairwise
nonisomorphism of those members.
\end{theorem}

\begin{proof}
\emph{Step 1: choose the local data before the pencil.}
Over $\mathbf Q$ choose $d$ hyperplanes in linear general position
and a degree-$d$ form $Q$ transverse to every positive-dimensional
intersection and avoiding the zero-dimensional intersections.
This is a nonempty open choice, by the ordinary Bertini theorem on
the finitely many projective linear strata.  Clear denominators to
obtain integral linear forms and an integral $Q$, and put
$P=\prod_i\ell_i$.  These geometric conditions persist modulo
all but finitely many primes by openness of smoothness and
properness of the relevant finite list of projective schemes.
Choose a prime $q>d$ outside this finite set, outside the prescribed
finite set, and with $q\nmid2d$.

Let
\[
 J=\sum_{i=0}^4X_i^{d-1}X_{i+1}+X_5^d,
 \qquad b_d=\prod_{p\mid d}p.
\]
Choose $F_1$ subject to
\[
       F_1\equiv P-dF_0+qQ\pmod{q^2},
       \qquad F_1\equiv J\pmod{b_d}.
 \tag{GF18.3}\label{r18gf:form-congruences}
\]
The moduli are coprime.  Their solution set is a full affine lattice
in the integral coefficient space and is Zariski dense over
$\mathbf Q$.  GP17's nonempty fixed-smooth-member Lefschetz open
therefore meets this lattice.  Take $F_1$ in that intersection.
The pencil $[u:v]\mapsto V(uF_0+vF_1)$ has exactly the geometric
properties of GP17; the singular special reduction at $q$ does
not alter this characteristic-zero application.

\emph{Step 2: spread only after fixing both members.}
By GP17 choose a finite set $S$ of rational primes containing
$q$ and every prime dividing $2d$, such that outside $S$ every
geometric member of the spread pencil is smooth or has only the
ordinary quadratic singularities allowed by that theorem.
The member $F_0$ is smooth modulo every $p\nmid d$, and the
congruence to $J$ makes $F_1$ smooth modulo every $p\mid d$.
These two fixed members supply the good reductions used at the
finitely many excluded primes.  Neither member is changed after
$S$ is chosen.

Use the Chinese remainder theorem to select the infinite progression
\[
 \begin{split}
 a&\equiv1\pmod{q^2},\qquad a\equiv1\pmod{b_d},\\
 a&\equiv0\pmod p
       \quad\text{for }p\in S,\quad p\nmid d,\quad p\ne q.
 \end{split}
 \tag{GF18.4}\label{r18gf:parameter-congruences}
\]
Discard the finitely many parameters giving singular
characteristic-zero pencil members.  The remaining $X_a$ are smooth
and geometrically integral by GP17's hypersurface integrality lemma.

\emph{Step 3: the exact local equation at $q$.}
For any chosen parameter define the integral form
\[
                         Q_a=\frac{dF_0+aF_1-P}{q}.
 \tag{GF18.5}\label{r18gf:Qa}
\]
The numerator is coefficientwise divisible by $q$ by
\eqref{r18gf:form-congruences} and
\eqref{r18gf:parameter-congruences}.  Moreover
$\bar Q_a=\bar Q$ modulo $q$, because $a\equiv1\pmod{q^2}$.
Thus over $\mathbf Z_q$ the exact equation is
\[
                              P+qQ_a=0.
 \tag{GF18.6}\label{r18gf:q-model}
\]
It has uniformizer exponent one, and the reduced hyperplane and
transversality data were fixed in Step 1.  Theorem~\ref{r18hp:model}
and Corollary~\ref{r18hp:U2} give full local ordinary actual G.
Theorem~\ref{r18hm:monodromy} gives the nonzero fourth monodromy.

\emph{Step 4: all other finite places.}
At $p\mid d$, equation \eqref{r18gf:family} reduces to the
smooth $F_1$.  At $p\in S$, $p\nmid d$, $p\ne q$, it reduces
to the smooth $F_0$.  These give proper smooth local models.
At $p\notin S$, the coefficient $d$ is a unit, and GP17's
spread-pencil statement and R17's all-prime ordinary-quadratic
activation prove full local G.  Possible node thickness is arbitrary
and is precisely within the scope of that frozen theorem.
No squarefreeness or boundedness condition on a discriminant value
is imposed.  This proves the first two assertions.

\emph{Step 5: all new cycles over the stated number fields.}
Let $F/\mathbf Q$ be unramified at $q$ and let $w\mid q$.
Then $q$ is still a uniformizer of $F_w$, up to a unit, and
\eqref{r18gf:q-model} has the same reduced data after its residue
field extension.  Apply Corollary~\ref{r18hp:U2} anew over $F_w$
to an arbitrary specified homologically trivial cycle there.
There is no descent assumption on this cycle.  At every other
place, good reduction remains good under scalar extension;
the ordinary-quadratic cases are handled anew by R17's theorem,
which allows the changed finite thickness and full native splitting
extensions.  Its finite auxiliary descents retain the complete
field degree and the original absolute value, exactly as in GP17.
Thus the full local groups, not merely pullbacks of earlier groups,
are covered at every finite place of $F$.

Apply the frozen PH16 global-height theorem to these full local
statements in each complementary Chow degree.  GA17 supplies the
admissible model lifts on this already established domain.  This
proves \eqref{r18gf:height}.  Finally, for a preassigned number
field avoid the finitely many ramified primes when choosing $q$.
\end{proof}

\paragraph{The exact enlargement and its limits.}
The number field varieties of \eqref{r18gf:family} have nonzero
$N^4$ at $q$.  Hence they cannot be among the varieties supplied
by GP17's good-or-ordinary-quadratic-at-every-place construction,
and cannot admit its more general dominant ordinary-quadratic
witness mechanism at $q$, by NM17.  The result is therefore an
additional full-height family with genuinely deeper degeneration,
rather than a change of model for a previously good generic fiber.
It makes no assertion that every hypersurface belongs to this
family and no assertion of stability of the full cycle domain under
arbitrary ramified extensions at $q$ without a further theorem.
General G and the full unconditional height for arbitrary varieties
remain open.

\begin{corollary}[Full codimension-two height domains in all $n\ge4$]
\label{r18gf:all-dimensions}
For every $n\ge4$ and every $d\ge n+2$, there are a prime
$q$, an integral degree-$d$ form $F_1$ in $n+2$ variables,
and an infinite arithmetic progression of integers $a$,
with finitely many excluded values, for which the smooth
geometrically integral hypersurfaces
\[
 X_a=\{dF_0+aF_1=0\}\subset\mathbf P^{n+1}_{\mathbf Q},
 \qquad F_0=\sum_{i=0}^{n+1}X_i^d,
 \tag{GF18.7}\label{r18gf:general-family}
\]
have ordinary actual strong G for every specified
geometrically homologically trivial codimension-two
rational cycle at every finite place.  They have
$N^n\ne0$ at $q$.  For every number field $F$ unramified
at $q$, all new local cycles in codimension two are
covered anew, and PH16 supplies the height on the full
groups
\[
 \operatorname{CH}^2((X_a)_F)_{\mathbf Q,\mathrm{geom\,hom}}
 \ \times\
 \operatorname{CH}^{n-1}((X_a)_F)_{\mathbf Q,\mathrm{geom\,hom}}
                         \longrightarrow\mathbf R.
 \tag{GF18.8}\label{r18gf:general-height}
\]
For a prescribed number field $F$, the construction may
choose $q$ unramified in $F$.  For $n>4$ this corollary
does not claim full G in the other middle codimensions.
\end{corollary}

\begin{proof}
Run the proof of Theorem~\ref{r18gf:global} in
$\mathbf P^{n+1}$ with $d$ general-position hyperplanes,
a degree-$d$ form transverse to every positive-dimensional
linear stratum and avoiding all zero-dimensional strata,
and
\[
                 J=\sum_{i=0}^{n}X_i^{d-1}X_{i+1}+X_{n+1}^d.
\]
GP17 is already valid in these dimensions and degrees.
Its fixed-$F_0$ open, the two coefficient congruences
modulo $q^2$ and $b_d$, the finite spreading set and
the parameter progression are unchanged.  In particular
at $q$ the exact equation is $P+qQ_a$ with
$\bar Q_a=\bar Q$.  Lemma~\ref{r18hg:model} and
Corollary~\ref{r18hg:G2} apply to this actual native
model, giving all codimension-two cycles.  At every
other place GP17 supplies full G, hence this degree.
Corollary~\ref{r18hm:all-dimensions} supplies $N^n\ne0$.

After a number-field extension unramified at $q$, rerun
the native model theorem on every cycle at the places
above $q$; at all other places rerun the already frozen
good or ordinary-quadratic theorem, with its permitted
changed thickness.  The complementary codimension
$n-1$ is the already frozen full one-cycle sector at
every finite place.  PH16 therefore applies on the
complete displayed pair of homologically trivial
Chow groups.  Avoiding the discriminant of a preassigned
number field when choosing $q$ proves the last assertion.
\end{proof}

\begin{corollary}[Zariski density while retaining the deep local model]
\label{r18gf:density}
Fix the dimension, degree, prime $q$, hyperplane product $P$
and transverse form $Q$ used in either construction above.
Allow $F_1$ to run through the fixed coefficient lattice
\eqref{r18gf:form-congruences} inside GP17's nonempty
Lefschetz open, and for each such $F_1$ run through the
permitted parameter progression chosen after spreading
that pencil.  The union of the resulting defining forms
$dF_0+aF_1$ is Zariski dense in the affine space of
degree-$d$ forms.  Hence their projective coefficient
classes are Zariski dense as well.
For $n=4$ every member has the full height domains of
Theorem~\ref{r18gf:global}; for $n\ge4$ every member has
the full displayed codimension-two height pairing.
The nonzero $N^n$ at the fixed $q$ is retained throughout.
\end{corollary}

\begin{proof}
Let $g$ be a rational polynomial in the coefficient
coordinates vanishing on this union.  For each admissible
$F_1$, the polynomial in one variable
$g(dF_0+tF_1)$ vanishes at infinitely many allowed integers
$t=a$.  It therefore vanishes identically.  Evaluating
at $t=1$ shows that $g(dF_0+F_1)=0$ for every admissible
$F_1$.  The fixed full coefficient lattice is Zariski
dense, and its intersection with the nonempty Lefschetz
open is still Zariski dense.  Translation by $dF_0$
preserves this property, so $g=0$.

The progression and finite spreading set may depend on
$F_1$; the argument needs only the infinitely many
permitted integers for each one.  All forms still have
the same reduction $P+qQ$ modulo $q^2$, so the already
proved local model, monodromy and height assertions
apply without any new uniform spreading assertion.
\end{proof}
